\section{Simulation Settings}
\label{sec:appendix-simulations}

\subsection{Models, Targets and Protocol Families}

Unless stated otherwise, $Z$ is a standardised Gaussian trajectory on a
uniform grid. For lag $u$, the stationary correlations are OU,
$\rho(u)=e^{-u/\tau}$; Mat\'ern-3/2,
$\rho(u)=(1+2u/\tau)e^{-2u/\tau}$, scaled so that
$\int_0^\infty \rho(u)\dd u=\tau$; a weighted two-scale OU mixture; and the
damped-periodic kernel $e^{-u/6}\cos(2\pi u/3)$. A between-unit trait share
$\alpha$ gives
$K_{jk}=\alpha+(1-\alpha)\rho(|t_j-t_k|)$. The five targets in the design
comparison are the temporal mean, occupation above thresholds $0$ and $1.5$,
two-sided occupation outside $[-1.2,1.2]$, and the logistic aggregate
$g(z)=\{1+e^{-2z}\}^{-1}$. The calibration experiments use the target subsets
specified below.

\begingroup
\setlength{\intextsep}{5pt}
\begin{table}[H]
\centering\footnotesize
\setlength{\tabcolsep}{2pt}
\renewcommand{\arraystretch}{0.94}
\begin{tabular}{@{}>{\raggedright\arraybackslash}p{0.16\textwidth}
>{\raggedright\arraybackslash}p{0.36\textwidth}
>{\raggedright\arraybackslash}p{0.44\textwidth}@{}}
\toprule
Experiment & Latent and calibration model & Actions, family and repetitions \\
\midrule
Uniform error,
\cref{fig:calibration}(a)
& $T=20$, $p=128$, OU $\tau=1$, $\alpha=0$; known calibration-noise variance
$0.25$, action-noise variance $0.5$; mean and occupation-at-zero targets
& 12 equally spaced point actions, exactly four selected
($\numEstFamily$ protocols); $m\in\{25,50,100,250,500,1000\}$ and
$\numEstReplications$ replications. \\[2pt]

Selection regret,
\cref{fig:calibration}(b)
& $T=20$, $p=128$; trait--OU $(\alpha,\tau)=(0.2,2)$ and two-scale OU
$(\tau_1,\tau_2,w)=(0.5,6,0.5)$; noiseless calibration and action-noise
variance $0.5$; mean and occupation targets at thresholds zero and one
& Points at $\{1,3,\ldots,19\}$, width-three windows centred at 5 and 15,
and pre-specified three-replicate point packages at 2 and 18; exactly four of
14 actions; the same $m$ grid and replication count as panel (a). \\[2pt]

Nested protocol classes,
\cref{fig:calibration}(c,d)
& $T=20$, $p=64$; noiseless calibration; occupation-at-zero target;
action-noise variance $0.4$;
$K_{jk}=0.15+0.85\exp\{-|t_j-t_k|/[\tau(t_j)\tau(t_k)]^{1/2}\}$, with
$\tau(t)$ linear from $0.3$ to $3$
& Exactly four actions are selected. L1 compares contiguous with dispersed
layouts; L2 adds early, middle, late and full-horizon phases; L3 uses supports
on eight coarse bins; L4 adds 495 fine-grid four-point subsets. The cumulative,
deduplicated class sizes are \numResClassSizeOne, \numResClassSizeTwo,
\numResClassSizeThree{} and \numResClassSizeFour; the same $m$ grid and 30
replications. \\
\bottomrule
\end{tabular}
\caption{Calibration and nested protocol-class configurations.}
\label{tab:synthetic-calibration-config}
\end{table}
\endgroup

For the target-aware design comparison in \cref{tab:s5}, all settings use
$T=10$, $p=128$ and $\alpha=0$. The stationary OU and Mat\'ern-3/2 settings
use $\tau=1$ and action-noise variance $0.4$. The horizon-varying setting uses
\[
K_{jk}
=
\exp\{-|t_j-t_k|/[\tau(t_j)\tau(t_k)]^{1/2}\},
\qquad
\tau(t)=0.5+3.5t/10,
\]
with action-noise variance $0.15$. The recency-weighted OU setting uses
$\tau=1$, target weights proportional to $e^{-(T-t)/6}$ and action-noise
variance $0.4$. The heterogeneous-action setting uses OU scale $\tau=0.7$
and time-dependent action-noise variance $1+2t/10$.

The first four settings use 12 point actions at
$\{0.5,0.5+9/11,\ldots,9.5\}$ and select exactly four actions. The
heterogeneous family uses six centres $0.5+1.8k$, $k=0,\ldots,5$, widths
$\{0,1.5,3.5\}$, cost $1+0.15w$ and a cost budget of four. All feasible
protocols are enumerated to obtain the target-specific optimum. The linear
comparator uses the actual action noise, kernel quadrature its standard
noiseless criterion, and mutual information and integrated posterior variance
the same covariance and action catalogue. Greedy search starts from the empty
protocol; one-swap refinement takes the best feasible exchange until no
improvement remains or \numSwapRounds{} exchanges have been accepted.

\subsection{Numerical Evaluation of Nonlinear Targets}
\label{sec:appendix-numerical}

For one-sided threshold targets, $C_g$ is evaluated from the Plackett integral.
If $G_c(r)$ denotes the corresponding one-sided covariance transform, then the
two-sided occupation target uses
\begin{equation*}
C_{{\rm two},c}(r)=2\{G_c(r)+G_c(-r)\}.
\end{equation*}
The logistic transform is evaluated through a truncated Hermite expansion,
with coefficients computed by Gauss--Hermite quadrature. 
